\let\negmedspace\undefined
\let\negthickspace\undefined
\documentclass[journal]{IEEEtran}
\usepackage[a5paper, margin=10mm, onecolumn]{geometry}
%\usepackage{lmodern} % Ensure lmodern is loaded for pdflatex
\usepackage{tfrupee} % Include tfrupee package

\setlength{\headheight}{1cm} % Set the height of the header box
\setlength{\headsep}{0mm}     % Set the distance between the header box and the top of the text

\usepackage{gvv-book}
\usepackage{gvv}
\usepackage{cite}
\usepackage{amsmath,amssymb,amsfonts,amsthm}
\usepackage{algorithmic}
\usepackage{graphicx}
\usepackage{textcomp}
\usepackage{xcolor}
\usepackage{txfonts}
\usepackage{listings}
\usepackage{enumitem}
\usepackage{mathtools}
\usepackage{gensymb}
\usepackage{comment}
\usepackage[breaklinks=true]{hyperref}
\usepackage{tkz-euclide} 
\usepackage{listings}
% \usepackage{gvv}                                        
\def\inputGnumericTable{}                                 
\usepackage[latin1]{inputenc}                                
\usepackage{color}                                            
\usepackage{array}                                            
\usepackage{longtable}                                       
\usepackage{calc}                                             
\usepackage{multirow}                                         
\usepackage{hhline}                                           
\usepackage{ifthen}                                           
\usepackage{lscape}
\begin{document}

\bibliographystyle{IEEEtran}
\vspace{3cm}

\title{1.4.27}
\author{EE25BTECH11011 - Navya}
% \maketitle
% \newpage
% \bigskip
{\let\newpage\relax\maketitle}

\renewcommand{\thefigure}{\theenumi}
\renewcommand{\thetable}{\theenumi}
\setlength{\intextsep}{10pt} % Space between text and floats
\textbf{Question}:\\
If $\vec{a}$ and $\vec{b}$ are the position vectors of $A$ and $B$ respectively, find the position vector of a point $C$ in $BA$ produced such that $BC$ = $1.5$ $BA$\\

\textbf{Solution:} \\
Given the position vectors of $\vec{A}$ and  $\vec{B}$  are  $\vec{a}$ and  $\vec{b}$ \\
Let the position vector of point $\vec{C}$ be  $\vec{c}$ 

The vector {BC} is given by:
\begin{align}
\vec{BC}=   1.5{BA}
\end{align}
\begin{align}
 (\vec{c}-\vec{a}) : (\vec{b}-\vec{c}) = 0.5 : 1.5 = 1 : 3 
\end{align}
Therefore, $ C $ divides $ AB $ externally in the ratio $ 1:3 $\\

Writing the ratio as $ k:1 $ 
we have 
\begin{align}
k=3
\end{align}
By the section formula (external division), the position vector of $C$  is\\

\begin{center}
\begin{align}
\vec{c}= \frac{k\vec{b} - \vec{a}}{k-1}, \quad (k=3).
\end{align}
\bigskip
\begin{align}
\vec{c} = \frac{3\vec{b} - \vec{a}}{2}.
\end{align}

\begin{align}
\therefore \quad \vec{C} =  \tfrac{3}{2}\vec{b} - \tfrac{1}{2}\vec{a} .
\end{align}

The above equation can be written as:
\begin{align}
\mathbf{c} = (\mathbf{a} \quad \mathbf{b})
\begin{pmatrix}
\frac{3}{2} \\
-\frac{1}{2}
\end{pmatrix}
\end{align}

\end{center}
\begin{figure}[h!]
    \centering
    \includegraphics[width=0.6\textwidth]{figs/fig1.png}
    \caption{Point C dividing AB externally in ratio 3:1}
    \label{fig:division}
\end{figure}
\end{document}